\section{Vollständige Induktion}

\dfn{Prinzip der vollständigen Induktion}{
  Sei $A(n)$ eine Aussage, die von einer natürlichen Zahl
  $n \in \bbN_0$ abhängt, wobei
  $\bbN_0 := \bbN \cup \{0\}$.

  Um zu zeigen, dass $A(n)$ für alle $n \in \bbN_0$ gilt,
  genügt es, die folgenden beiden Schritte zu beweisen:

  \begin{enumerate}[label=\textbf{(\arabic*)}]
    \item \textbf{Induktionsanfang:} Die Aussage $A(0)$ ist wahr.

    \item \textbf{Induktionsschritt:} Für jedes $n \in \bbN_0$
    gilt die Implikation
    \[
      A(n) \implies A(n+1).
    \]
  \end{enumerate}

  Dann gilt $A(n)$ für alle $n \in \bbN_0$.
}

\nt{
  Die Idee der vollständigen Induktion ist, dass man eine Aussage
  zunächst für den Startwert beweist und anschließend zeigt, dass
  ihre Gültigkeit immer auf die nächste natürliche Zahl übertragen
  werden kann.

  Ist also $A(0)$ wahr, so folgt aus dem Induktionsschritt zunächst
  $A(1)$, daraus $A(2)$, daraus $A(3)$ und so weiter.

  Wichtig ist, dass man im Induktionsschritt nicht einfach annimmt,
  dass die Behauptung für alle Zahlen gilt. Man nimmt sie nur für
  ein beliebiges, aber festes $n$ an und beweist daraus die Aussage
  für $n+1$.
}

\ex{Summenformel}{
  Wir zeigen, dass für alle $n \in \bbN_0$ gilt:
  \[
    \sum_{k=0}^{n} k = \frac{n(n+1)}{2}.
  \]

  \begin{myproof}
    Wir führen den Beweis durch vollständige Induktion über $n$.

    \textbf{Induktionsanfang:}

    Für $n=0$ ist die linke Seite gleich $0$. Somit gilt
    \[
      \sum_{k=0}^{0} k
      = 0
      = \frac{0\cdot(0+1)}{2}.
    \]
    Die Behauptung gilt also für $n=0$.

    \textbf{Induktionsannahme:}

    Sei $n \in \bbN_0$ beliebig und es gelte
    \[
      \sum_{k=0}^{n} k = \frac{n(n+1)}{2}.
    \]

    \textbf{Induktionsschritt:}

    Wir zeigen die Behauptung für $n+1$. Dazu verwenden wir
    die Induktionsannahme:
    \begin{align*}
      \sum_{k=0}^{n+1} k
      &= \sum_{k=0}^{n} k + (n+1) \\
      &= \frac{n(n+1)}{2} + (n+1) \\
      &= \frac{n(n+1)+2(n+1)}{2} \\
      &= \frac{(n+1)(n+2)}{2}.
    \end{align*}

    Dies ist genau die behauptete Formel mit $n+1$ anstelle
    von $n$. Damit ist der Induktionsschritt bewiesen.

    Folglich gilt
    \[
      \sum_{k=0}^{n} k = \frac{n(n+1)}{2}
      \qquad\text{für alle } n \in \bbN_0.
      \qedhere
    \]
  \end{myproof}
}

\nt{
  Die Summenformel lässt sich auch durch einen Trick zeigen,
  der häufig Gauß zugeschrieben wird. Man schreibt die Summe
  einmal vorwärts und einmal rückwärts:
  \begin{align*}
    S &= 0+1+2+\dots+n,\\
    S &= n+(n-1)+(n-2)+\dots+0.
  \end{align*}
  Addiert man beide Zeilen, erhält man
  \[
    2S = n(n+1).
  \]
  Also folgt $S=\frac{n(n+1)}{2}$.
}

\ex{Teilbarkeit}{
  Wir zeigen, dass
  \[
    7 \mid \bigl(2^{3n}-1\bigr)
    \qquad\text{für alle } n\in\bbN_0.
  \]

  \begin{myproof}
    Wir führen eine vollständige Induktion über $n$.

    \textbf{Induktionsanfang:}

    Für $n=0$ gilt
    \[
      2^{3\cdot 0}-1 = 1-1=0.
    \]
    Da $0=7\cdot 0$ ist, gilt $7\mid 0$.

    \textbf{Induktionsannahme:}

    Sei $n\in\bbN_0$ und es gelte
    \[
      7\mid \bigl(2^{3n}-1\bigr).
    \]
    Das bedeutet, dass es ein $q\in\bbZ$ gibt mit
    \[
      2^{3n}-1=7q.
    \]

    \textbf{Induktionsschritt:}

    Wir betrachten den Ausdruck für $n+1$:
    \begin{align*}
      2^{3(n+1)}-1
      &= 2^3\cdot 2^{3n}-1\\
      &= 8\cdot 2^{3n}-1\\
      &= 7\cdot 2^{3n}+(2^{3n}-1).
    \end{align*}
    Nach der Induktionsannahme ist $2^{3n}-1$ durch $7$
    teilbar. Auch $7\cdot 2^{3n}$ ist durch $7$ teilbar.
    Daher ist ihre Summe durch $7$ teilbar.

    Somit gilt
    \[
      7\mid \bigl(2^{3(n+1)}-1\bigr).
    \]
    Der Induktionsschritt ist bewiesen, also gilt die Behauptung
    für alle $n\in\bbN_0$.
  \end{myproof}
}

\dfn{Induktion ab einem Startwert}{
  Sei $k\in\bbN_0$ und $A(n)$ eine Aussage für $n\geq k$.
  Gilt
  \begin{enumerate}[label=\textbf{(\arabic*)}]
    \item $A(k)$,
    \item $A(n)\implies A(n+1)$ für alle $n\geq k$,
  \end{enumerate}
  so gilt $A(n)$ für alle $n\geq k$.
}

\dfn{Starke Induktion}{
  Sei $k\in\bbN_0$ und $A(n)$ eine Aussage für $n\geq k$.
  Gilt $A(k)$ und lässt sich für jedes $n\geq k$ zeigen, dass
  \[
    \bigl(A(k)\land A(k+1)\land\dots\land A(n)\bigr)
    \implies A(n+1),
  \]
  so gilt $A(n)$ für alle $n\geq k$.
}

\nt{
  Bei der gewöhnlichen Induktion verwendet man im Induktionsschritt
  nur die Aussage $A(n)$. Bei der starken Induktion darf man
  dagegen alle vorherigen Aussagen
  $A(k),A(k+1),\dots,A(n)$ verwenden.

  Beide Varianten beruhen auf dem Prinzip der vollständigen
  Induktion. Welche Variante geeigneter ist, hängt davon ab,
  welche Informationen man für den nächsten Schritt benötigt.
}

